% ============================================================================
\section{Example II: a driven two-level system}
\label{sec:driven}
% ============================================================================

\subsection{The lab-frame Hamiltonian}

A spin-$1/2$ nucleus in a static field $B_0\hat{z}$ has the Zeeman Hamiltonian
%
\begin{equation}
\Hop_0 = -\gamma B_0 \Szo = \frac{\hbar\omega_0}{2}\szo,
\qquad \omega_0 \equiv -\gamma B_0 ,
\end{equation}
%
pure $\szo$, with the two spin states split by the Larmor energy $\hbar\omega_0$. For a
proton at \SI{11.7}{\tesla} this is \SI{500}{\mega\hertz}, the number printed on the
front of the spectrometer. Nothing happens under $\Hop_0$ alone: $\ket{\alpha}$ and
$\ket{\beta}$ are eigenstates and the Bloch vector merely spins about $\hat{z}$.

To make something happen, add a weak field rotating in the $xy$ plane at frequency
$\omega$ --- the RF coil in NMR, or, for an electronic transition in the dipole
approximation, the optical field of a laser:
%
\begin{equation}
\Hop(t) = \frac{\hbar\omega_0}{2}\szo
  + \frac{\hbar\omega_1}{2}\Big[\cos(\omega t+\varphi)\,\sxo + \sin(\omega t+\varphi)\,\syo\Big],
\qquad \omega_1 \equiv -\gamma B_1 .
\label{eq:lab-frame}
\end{equation}
%
This is genuinely time dependent, so the machinery of \cref{sec:geometry} does not apply
directly. One transformation fixes that.

\subsection{The rotating frame}

Move to a frame rotating about $z$ at the drive frequency by defining
$\ket{\tilde\psi} = \Ry(t)\ket{\psi}$ with $\Ry(t) = e^{+\ii\omega t\szo/2}$. Substituting
into the Schr\"odinger equation (worked out step by step in \cref{app:rotframe}) gives
an effective Hamiltonian $\tilde{H} = \Ry\Hop\Ry^\dagger - \tfrac{\hbar\omega}{2}\szo$
which turns out to be \emph{time independent}:

\begin{keyresult}[Rotating-frame Hamiltonian]
\vspace{-0.35cm}
\begin{equation}
\Hop_{\rm rot} = \frac{\hbar}{2}\Big[\underbrace{(\omega_0-\omega)}_{\textstyle \Delta}\,\szo
   + \omega_1\big(\cos\varphi\,\sxo + \sin\varphi\,\syo\big)\Big],
\label{eq:rotframe}
\end{equation}
that is, $d_z = \hbar\Delta$, $d_x = \hbar\omega_1\cos\varphi$,
$d_y = \hbar\omega_1\sin\varphi$.
\end{keyresult}

\noindent
Three features of \cref{eq:rotframe} are each a piece of real spectroscopic practice.

\begin{itemize}[leftmargin=*,itemsep=3pt]
\item \textbf{The detuning is the $\szo$ knob.} On resonance ($\omega=\omega_0$) it
vanishes and $\dvec$ lies flat in the equatorial plane. Off resonance, $\dvec$ tilts
toward the poles --- and by \cref{fig:bloch-prec}, a tilted $\dvec$ means the state can
no longer reach the far pole.

\item \textbf{The pulse phase is the $\syo$ knob.} Setting $\varphi=0$ gives a pure
$\sxo$ term, an ``$x$-pulse''; $\varphi=\pi/2$ gives a pure $\syo$ term, a ``$y$-pulse''.
This is not jargon: $\varphi$ is literally the phase of the RF pulse, and it is why every
pulse-sequence diagram in an NMR textbook labels its pulses $x$, $y$, $-x$, $-y$.

\item \textbf{The tilted-cone geometry of \cref{sec:bloch} is now physical.} The polar
angle of $\dvec$ is set by how far off resonance you are, and its azimuth by the pulse
phase.
\end{itemize}

\subsection{Rabi oscillations}

Apply \cref{eq:eigenvalues}. The splitting in the rotating frame defines the
\textbf{generalized Rabi frequency}
%
\begin{equation}
\Omega = \frac{|\dvec|}{\hbar} = \sqrt{\Delta^2+\omega_1^2},
\label{eq:rabi-freq}
\end{equation}
%
and integrating \cref{eq:bloch-eq} from $\svec(0)=\hat{z}$ (or equivalently applying
\cref{eq:eigenvectors} twice) gives

\begin{keyresult}[Rabi formula]
\vspace{-0.35cm}
\begin{equation}
P_\beta(t) = \frac{\omega_1^2}{\Omega^2}\,\sin^2\!\left(\frac{\Omega t}{2}\right).
\label{eq:rabi}
\end{equation}
\end{keyresult}

\noindent
\Cref{eq:rabi} repays careful reading. The \emph{rate} of oscillation,
$\Omega=\sqrt{\Delta^2+\omega_1^2}$, always \emph{increases} with detuning. But the
\emph{amplitude},
%
\begin{equation}
P_\beta^{\max} = \frac{\omega_1^2}{\Delta^2+\omega_1^2},
\label{eq:lorentzian}
\end{equation}
%
is a Lorentzian in $\Delta$ of half-width $\omega_1$: complete population inversion is
possible only exactly on resonance. \Cref{eq:lorentzian} is simultaneously the origin of
the spectroscopic lineshape and of \emph{pulse selectivity} --- a weak, long pulse
addresses a narrow band of frequencies; a strong, short pulse excites everything in
sight. \Cref{fig:rabi} shows both effects.

% ---------------------------------------------------------------------------
\begin{figure}[t]
\centering
\begin{tikzpicture}
\begin{groupplot}[
  group style={group size=2 by 1, horizontal sep=1.5cm},
  notestyle, height=6.0cm,
]
\nextgroupplot[
  xlabel={$\omega_1 t$}, ylabel={$P_\beta(t)$},
  title={Rabi oscillations},
  xmin=0, xmax=12.566, ymin=-0.03, ymax=1.05,
  xtick={0,3.1416,6.2832,9.4248,12.566},
  xticklabels={$0$,$\pi$,$2\pi$,$3\pi$,$4\pi$},
  legend style={at={(0.5,-0.30)},anchor=north,legend columns=4,draw=none,fill=none,
                font=\footnotesize},
  legend cell align=left,
]
\addplot[cLower, very thick, domain=0:12.566, samples=400]
  {(1/1)*sin(deg(sqrt(1)*x/2))^2};                       \addlegendentry{$\Delta=0$}
\addplot[cAccent, very thick, domain=0:12.566, samples=400]
  {(1/2)*sin(deg(sqrt(2)*x/2))^2};                       \addlegendentry{$\Delta=\omega_1$}
\addplot[cWarm, very thick, domain=0:12.566, samples=500]
  {(1/5)*sin(deg(sqrt(5)*x/2))^2};                       \addlegendentry{$\Delta=2\omega_1$}
\addplot[cUpper, very thick, domain=0:12.566, samples=600]
  {(1/17)*sin(deg(sqrt(17)*x/2))^2};                     \addlegendentry{$\Delta=4\omega_1$}

\nextgroupplot[
  xlabel={detuning $\Delta/\omega_1$}, ylabel={$P_\beta^{\max}$},
  title={\ldots but only on resonance is inversion complete},
  xmin=-5, xmax=5, ymin=-0.03, ymax=1.06,
]
\addplot[cUpper, very thick, domain=-5:5, samples=300]{1/(1+x^2)};
\addplot[cMuted, dotted, thin, domain=-5:5, samples=2]{0.5};
\draw[<->,cAccent,thin] (axis cs:-1,0.5) -- (axis cs:1,0.5);
\node[cAccent,font=\footnotesize,anchor=south] at (axis cs:0,0.53){FWHM $=2\omega_1$};
\end{groupplot}
\end{tikzpicture}
\caption{\textbf{Left:} \cref{eq:rabi} at four detunings. Detuning makes the
oscillation \emph{faster} and \emph{shallower} at the same time. \textbf{Right:} the
maximum transferable population, \cref{eq:lorentzian}. A pulse of Rabi strength
$\omega_1$ can only fully invert spins within roughly $\pm\omega_1$ of resonance ---
which is what ``selective pulse'' means.}
\label{fig:rabi}
\end{figure}
% ---------------------------------------------------------------------------

\subsection{$\pi$ and $\pi/2$ pulses}

On resonance $\dvec$ lies in the equatorial plane, the precession cone opens into a
great circle through both poles, and the state sweeps from $\ket{\alpha}$ to
$\ket{\beta}$ and back. Two durations matter:
%
\begin{equation}
t_\pi = \frac{\pi}{\omega_1}: \;\; \ket{\alpha}\to\ket{\beta} \text{ (inversion)},
\qquad
t_{\pi/2} = \frac{\pi}{2\omega_1}: \;\; \ket{\alpha}\to \text{the equator (coherence)}.
\end{equation}
%
The $\pi/2$ pulse is the one that generates an NMR signal: it converts population
difference into transverse coherence, which is what the detection coil sees as the free
induction decay. And the pulse phase decides \emph{where} on the equator you land. Using
\cref{eq:bloch-eq} with $\dvec = \hbar\omega_1\hat{x}$,
%
\begin{equation}
\text{$x$-pulse: } \hat{z}\to-\hat{y},
\qquad
\text{$y$-pulse: } \hat{z}\to+\hat{x},
\end{equation}
%
two states $\ang{90}$ apart on the equator. Both are ``the same pulse'' in every respect
except phase.

\begin{lookahead}[Where this is going: the Jaynes--Cummings model]
Write the transverse part of \cref{eq:rotframe} using the ladder operators
$\hat\sigma_\pm = \half(\sxo\pm\ii\syo)$:
\begin{equation}
\frac{\hbar\omega_1}{2}\big(\cos\varphi\,\sxo+\sin\varphi\,\syo\big)
 = \frac{\hbar\omega_1}{2}\Big(e^{-\ii\varphi}\hat\sigma_+ + e^{+\ii\varphi}\hat\sigma_-\Big).
\end{equation}
The drive \emph{raises} the two-level system while supplying a phase, or \emph{lowers} it
while absorbing one. Here the field is a classical number: $\omega_1\propto B_1$, and it
is inexhaustible --- the system can cycle forever and the field never notices.

Now promote the field to a quantum degree of freedom. The phase factors become photon
annihilation and creation operators, $e^{-\ii\varphi}\to\hat{a}$ and
$e^{+\ii\varphi}\to\hat{a}^\dagger$, and
\begin{equation}
\Hop_{\rm JC} = \frac{\hbar\omega_0}{2}\szo + \hbar\omega\,\hat{a}^\dagger\hat{a}
  + \hbar g\big(\hat{a}\,\hat\sigma_+ + \hat{a}^\dagger\hat\sigma_-\big),
\label{eq:jc}
\end{equation}
the \textbf{Jaynes--Cummings} Hamiltonian \cite{JaynesCummings} --- the single-emitter
case of the Tavis--Cummings model built in Computational Set 4. The Hilbert space is no
longer two dimensional; it is spin $\otimes$ boson. The classical drive strength
$\omega_1$ is replaced by $2g\sqrt{n+1}$, so that even the vacuum ($n=0$) drives Rabi
oscillations. That is vacuum Rabi splitting, and it is the foundation of polaritonic
chemistry. Everything in this section is its classical-field limit.
\end{lookahead}

\begin{nbbox}[In the notebook]
Part 5 verifies \cref{eq:rabi} against numerical propagation to $10^{-9}$ for several
detunings, confirms that a resonant $\pi$ pulse inverts the population exactly, and
shows that the $x$- and $y$-pulses leave the Bloch vector $\ang{90}$ apart.
\end{nbbox}
